\documentclass[10pt,letterpaper]{article}
\usepackage[T1]{fontenc}
\usepackage{amsmath,amssymb,amsthm}
\usepackage[margin=1in]{geometry}
\usepackage[most]{tcolorbox}
\usepackage{microtype}
\usepackage[hidelinks]{hyperref}
\usepackage{needspace}
\tcbset{colback=black!4,colframe=black!75,boxrule=0.7pt,arc=2pt,
 left=10pt,right=10pt,top=7pt,bottom=7pt}
\setlength{\parindent}{0pt}
\setlength{\parskip}{5pt}
\setlength{\emergencystretch}{2em}
\title{Angle Chasing}
\author{Adhiraj Singh Anand}
\date{November 7, 2025}

\begin{document}\maketitle\section{Examples}
\Needspace{5\baselineskip}
\begin{tcolorbox}[breakable,title={IMO 2018 P1},colbacktitle=black!4,coltitle=black,fonttitle=\normalfont]
\textbf{Example 1:} Let $\Gamma$ be the circumcircle of acute triangle $ABC$. Points $D$ and $E$ are on segments $AB$ and $AC$ respectively such that $AD = AE$. The perpendicular bisectors of $BD$ and $CE$ intersect minor arcs $AB$ and $AC$ of $\Gamma$ at points $F$ and $G$ respectively. Prove that lines $DE$ and $FG$ are either parallel or they are the same line.

\end{tcolorbox}
\textit{Proof:}
Extend $FD, GE$ to $X,Y$ on $\Gamma$.

$\angle AXF=\angle ABF=\angle FDB=\angle ADX$ so $AD=AX=AY$.

So, $Y,D,E,X$ are concyclic with circumcenter $A$.

Now, $\angle YED=\angle YXD=\angle YXF=\angle YGF$.

Corresponding angles are equal, so $DE\parallel FG\square$.
\nopagebreak[4]\par\vspace{5pt}{\footnotesize\href{https://artofproblemsolving.com/community/u1205624h1670580p34930519}{AoPS solution: me}}
\par\vspace{12pt}
\Needspace{5\baselineskip}
\begin{tcolorbox}[breakable,title={Indian TST 2019 Practice Test 2 P1},colbacktitle=black!4,coltitle=black,fonttitle=\normalfont]
\textbf{Example 2:} Let the points $O$ and $H$ be the circumcenter and orthocenter of an acute angled triangle $ABC.$ Let $D$ be the midpoint of $BC.$ Let $E$ be the point on the angle bisector of $\angle BAC$ such that $AE\perp HE.$ Let $F$ be the point such that $AEHF$ is a rectangle. Prove that $D,E,F$ are collinear.
\end{tcolorbox}
\textit{Proof:}
Let $XYZ$ be the orthic triangle, $AF\cap YZ=U$, and $EF \cap YZ=V$.

Notice that $A,H,E,F,Y,Z$ are concyclic with centre $U$, and $B,C,Y,Z$ are concyclic with centre $D$. Additionally, \[\angle ZVE=180^{\circ}-\angle YZE-\angle FEZ=180^{\circ}-\angle ZAE-\angle FAZ=90^{\circ}\]Thus, $D,E,F,$ and $G$ all lie on the perpendicular bisector of the radical axis $YZ$, so we're done. $\square$
\nopagebreak[4]\par\vspace{5pt}{\footnotesize\href{https://artofproblemsolving.com/community/u1205624h1876805p36381082}{AoPS solution: me}}
\par\vspace{12pt}
\Needspace{5\baselineskip}
\begin{tcolorbox}[breakable,title={IMO Shortlist 2020 G3},colbacktitle=black!4,coltitle=black,fonttitle=\normalfont]
\textbf{Example 3:} Let $ABCD$ be a convex quadrilateral with $\angle ABC>90$, $CDA>90$ and $\angle DAB=\angle BCD$. Denote by $E$ and $F$ the reflections of $A$ in lines $BC$ and $CD$, respectively. Suppose that the segments $AE$ and $AF$ meet the line $BD$ at $K$ and $L$, respectively. Prove that the circumcircles of triangles $BEK$ and $DFL$ are tangent to each other.
\end{tcolorbox}
\textit{Proof:}
Let $X$ be the reflection of $A$ over $BD$.

Notice $\angle DA'L=\angle DAL=\angle DFA$, so $A'$ lies on $(DFL)$.

Similarly, $\angle BA'K=\angle BAK=\angle BEK$, so $A'$ lies on $(BEK)$.

Now, notice $\measuredangle A'LD+\measuredangle A'KE=\measuredangle A'LK+\measuredangle A'KL=\measuredangle KA'L$, so we are done.
\nopagebreak[4]\par\vspace{5pt}{\footnotesize\href{https://artofproblemsolving.com/community/u1205624h2625881p35031833}{AoPS solution: me}}
\par\vspace{12pt}
\Needspace{5\baselineskip}
\begin{tcolorbox}[breakable,title={APMO 2024 P5},colbacktitle=black!4,coltitle=black,fonttitle=\normalfont]
\textbf{Example 4:} Line $\ell$ intersects sides $BC$ and $AD$ of cyclic quadrilateral $ABCD$ in its interior points $R$ and $S$, respectively, and intersects ray $DC$ beyond point $C$ at $Q$, and ray $BA$ beyond point $A$ at $P$. Circumcircles of the triangles $QCR$ and $QDS$ intersect at $N \neq Q$, while circumcircles of the triangles $PAS$ and $PBR$ intersect at $M\neq P$. Let lines $MP$ and $NQ$ meet at point $X$, lines $AB$ and $CD$ meet at point $K$ and lines $BC$ and $AD$ meet at point $L$. Prove that point $X$ lies on line $KL$.
\end{tcolorbox}
\textit{Proof:}
Notice that $M$ is the Miquel point of complete quadrilateral $ABRSPL$, and $N$ of $CDSRQL$.

So, $M,N,L,R,S$ are concyclic.

\begin{align*}\measuredangle PMN&=\measuredangle PMS+\measuredangle SMN\\&=\measuredangle PAS+\measuredangle SLN\\&=\measuredangle BAS+\measuredangle SRN\\&=\measuredangle RCD+\measuredangle QRN\\&=\measuredangle RNQ+\measuredangle QRN=\measuredangle PQN\end{align*}Thus, $P,M,N,Q$ are concyclic. Let $PX\cap (MSRN)=Y$, $QX\cap (MSRN)=Z$.

By Reim's theorem $\triangle LYZ$ and $\triangle KPQ$ are homothetic with centre $X$, so $K,X,L$ are collinear as desired. $\square$
\nopagebreak[4]\par\vspace{5pt}{\footnotesize\href{https://artofproblemsolving.com/community/u1205624h3368255p35157277}{AoPS solution: me}}
\par\vspace{12pt}\newpage
\section{Practise Problems}
Here are some practise problems, in increasing order of difficulty.
\begin{enumerate}
\item \textbf{\href{https://www.imo-official.org/assets/documents/problems/2007/2007_eng.pdf}{IMO 2007 P4}}

In triangle $ABC$ the bisector of angle $BCA$ intersects the circumcircle again at $R$, the perpendicular bisector of $BC$ at $P$, and the perpendicular bisector of $AC$ at $Q$. The midpoint of $BC$ is $K$ and the midpoint of $AC$ is $L$. Prove that the triangles $RPK$ and $RQL$ have the same area.
\item \textbf{\href{https://www.imo-official.org/assets/documents/problems/2010/2010_eng.pdf}{IMO 2010 P4}}

Let $P$ be a point inside the triangle $ABC$. The lines $AP$, $BP$ and $CP$ intersect the circumcircle $\Gamma$ of triangle $ABC$ again at the points $K$, $L$ and $M$ respectively. The tangent to $\Gamma$ at $C$ intersects the line $AB$ at $S$. Suppose that $SC=SP$. Prove that $MK=ML$.
\item \textbf{\href{https://www.imo-official.org/assets/documents/problems/2015/2015_eng.pdf}{IMO 2015 P3}}

Let $ABC$ be an acute triangle with $AB>AC$. Let $\Gamma$ be its circumcircle, $H$ its orthocentre, and $F$ the foot of the altitude from $A$. Let $M$ be the midpoint of $BC$. Let $Q$ be the point on $\Gamma$ such that $\angle HQA=90^\circ$, and let $K$ be the point on $\Gamma$ such that $\angle HKQ=90^\circ$. Assume that the points $A,B,C,K$ and $Q$ are all different, and lie on $\Gamma$ in this order. Prove that the circumcircles of triangles $KQH$ and $FKM$ are tangent to each other.
\end{enumerate}\end{document}